\begin{theorem}[Project-local one-site zero convention]
    \label{thm:project_local_one_site_parent_hamiltonian_zero}
    \lean{MPSTensor.parentHamiltonian_two_one_eq_zero}
    \lean{MPSTensor.parentHamiltonianGroundSpaceES_two_one_eq_top}
    \leanok
    For the convention that an $L$-site interaction is zero when $L>N$, the
    two-site parent Hamiltonian on a one-site periodic chain satisfies
    \begin{align}
        H_2^{(1)}      & = 0, \notag    \\
        \ker H_2^{(1)} & = \C^d. \notag
    \end{align}
    This convention is local to the present development. The two-site terms
    in \cite[Section~III]{Beigi2012Classification} are
    not assigned an action when both site labels denote the same site. The
    nearest-neighbor condition in
    \cite[Theorem~3.10(iii)]{Cirac2016MPDO_arXiv}, source lines 534--540, only
    concerns $N>2$.
\end{theorem}

\begin{proof}\leanok
    Since $2>1$, the single translated interaction is zero by the stated
    convention. Its sum is therefore zero, and the kernel of the zero
    operator is the full one-site Hilbert space.
\end{proof}

\begin{theorem}[Classification of eventually constant weighted sector graphs]
    \label{thm:beigi_eventually_constant_sector_graph}
    \lean{FiniteWeightedDigraph.cyclicEdgeWeight_eq_one}
    \lean{FiniteWeightedDigraph.cycle_eq_constant}
    \lean{FiniteWeightedDigraph.cycleWeightSum_eq_card_loop}
    \leanok
    \uses{def:beigi_sector_graph}
    Let $G$ be a finite directed graph whose edge weights are non-negative
    integers, and let
    \begin{align}
        s_N
         & =
        \sum_{(a_0,\ldots,a_{N-1})}
        \prod_{n=0}^{N-1}w(a_n,a_{n+1}) \notag
    \end{align}
    be the weighted sum over ordered cyclic vertex sequences. Suppose that
    $s_N$ is constant for all sufficiently large $N$. Then every edge that
    belongs to a directed cycle has weight one, every directed cycle is a
    repetition of a single loop, and, for every $N>0$,
    $s_N=\#\{a:w(a,a)>0\}$. This is the comparison in
    \cite[Section~IV]{Beigi2012Classification}.
\end{theorem}

\begin{proof}\leanok
    Choose a sufficiently large $N$ such that the sums at lengths $N$,
    $N+1$, and $N(N+1)$ have their common eventual value. Repeating each
    ordered $N$-cycle $N+1$ times gives an injection into the cycles of length
    $N(N+1)$, while repeating each ordered $(N+1)$-cycle $N$ times gives a
    second injection. Comparison of the three weighted sums shows first that
    every positive cyclic weight is one. Both injections must then be
    exhaustive. Hence a cycle at the common length has periods $N$ and
    $N+1$; comparison of adjacent entries gives period one. Repeating an
    arbitrary shorter cycle to a sufficiently large length proves that it is
    constant. The only surviving ordered cycles at any positive length are
    therefore the constant words supported on positive loops, each of weight
    one.
\end{proof}

\begin{corollary}[Parent-ground-space dimension is the loop number]
    \label{cor:beigi_parent_ground_space_loop_number}
    \lean{MPSTensor.BeigiSectorGraphData.parentHamiltonianGroundSpaceES_finrank_eq_card_loop}
    \lean{MPSTensor.BeigiSectorGraphData.parentHamiltonianGroundSpaceES_finrank_eq_card_loop_of_hasBNTSectorData}
    \leanok
    \uses{thm:beigi_ordered_cycle_ground_space_dimension,
        thm:beigi_eventually_constant_sector_graph,
        thm:nncph_eventual_ground_space_dimension}
    If the parent-ground-space dimension is eventually constant, then for
    every $N>2$ it equals the number of positive loops in the sector graph.
    In particular, this conclusion holds under the BNT and all-chain
    nearest-neighbor parent-ground-space hypotheses of
    Theorem~\ref{thm:nncph_eventual_ground_space_dimension}.
\end{corollary}

\begin{proof}\leanok
    The ordered-cycle dimension formula identifies the eventual
    parent-ground-space dimension with the eventual weighted-cycle sum.
    The preceding theorem identifies that sum, at every positive length, with
    the loop number. Transport between coefficient and Euclidean coordinates
    preserves the dimension of the ground space.
\end{proof}

\begin{definition}[Tensor associated with a positive Beigi loop]
    \label{def:beigi_loop_product_state}
    \lean{MPSTensor.BeigiSectorGraphData.loopBondVector}
    \lean{MPSTensor.BeigiSectorGraphData.loopBondVector_mem}
    \lean{MPSTensor.BeigiSectorGraphData.loopBondVector_ne_zero}
    \lean{MPSTensor.BeigiSectorGraphData.loopCoordinateTensor}
    \lean{MPSTensor.BeigiSectorGraphData.loopTensor}
    \lean{MPSTensor.BeigiSectorGraphData.loopCyclicProduct}
    \lean{MPSTensor.BeigiSectorGraphData.loopProductState}
    \lean{MPSTensor.BeigiSectorGraphData.loopProductStateES}
    \leanok
    \uses{def:beigi_sector_graph}
    Let $a\to a$ be a positive loop and choose a nonzero vector
    $\varphi_a\in\ran P_{a,a}$. In sector coordinates define the matrix
    product tensor $(C_a^{(r,l)})_{x,y}=\delta_{x,l}\varphi_a(r,y)$ and set it
    to zero outside sector~$a$. The physical loop tensor is the
    image of $C_a$ under the one-site unitary of the spatial decomposition.
    Its associated cyclic product pairs the right factor at site $n$ with the
    left factor at site $n+1$. The corresponding $N$-site product state is
    denoted by $\Psi_a^{(N)}\in(\C^d)^{\otimes N}$.
\end{definition}

\begin{definition}[Schmidt-support tensor associated with a positive Beigi loop]
    \label{def:beigi_loop_schmidt_support}
    \lean{MPSTensor.BeigiSectorGraphData.loopSchmidtRank}
    \lean{MPSTensor.BeigiSectorGraphData.LoopSchmidtFactorization}
    \lean{MPSTensor.BeigiSectorGraphData.loopSchmidtFactorization}
    \lean{MPSTensor.BeigiSectorGraphData.minimalLoopCoordinateTensor}
    \lean{MPSTensor.BeigiSectorGraphData.minimalLoopTensor}
    \lean{MPSTensor.BeigiSectorGraphData.minimalLoopCoordinateTensor_sector_apply}
    \lean{MPSTensor.BeigiSectorGraphData.minimalLoopCoordinateTensor_sector_ne}
    \leanok
    \uses{def:beigi_loop_product_state}
    Let $\Phi_a=(\varphi_a(r,l))_{r,l}$ and write
    $r_a=\rank\Phi_a$. Choose a rank factorization
    \begin{align}
        \Phi_a & = X_aY_a, \notag                        \\
        X_a    & \in M_{D_a^{\mathrm R},r_a}(\C), \notag \\
        Y_a    & \in M_{r_a,D_a^{\mathrm L}}(\C). \notag
    \end{align}
    The tensor on the minimal Schmidt support is
    \begin{align}
        (\widetilde C_a^{(r,l)})_{\alpha,\beta}
         & =
        (Y_a)_{\alpha,l}(X_a)_{r,\beta}, \notag
    \end{align}
    again extended by zero outside sector~$a$. Its physical tensor is its
    image under the same one-site unitary.
\end{definition}

\begin{definition}[Normalized minimal tensor associated with a positive Beigi loop]
    \label{def:beigi_normalized_minimal_loop_tensor}
    \lean{MPSTensor.BeigiSectorGraphData.loopBondNorm}
    \lean{MPSTensor.BeigiSectorGraphData.normalizedMinimalLoopCoordinateTensor}
    \lean{MPSTensor.BeigiSectorGraphData.normalizedMinimalLoopTensor}
    \leanok
    \uses{def:beigi_loop_schmidt_support}
    Let $\nu_a=\lVert\varphi_a\rVert_2>0$, where the norm is the
    Hilbert-space norm of the bond vector, or equivalently the Frobenius norm
    of $\Phi_a$. In either coordinate system, the normalized representative is
    $\widehat C_a=\nu_a^{-1}\widetilde C_a$. The raw coordinate and physical
    tensors are retained separately.
\end{definition}

\begin{lemma}[Positive Schmidt dimension of a Beigi loop]
    \label{lem:beigi_loop_schmidt_rank_pos}
    \lean{MPSTensor.BeigiSectorGraphData.loopSchmidtRank_pos}
    \leanok
    \uses{def:beigi_loop_schmidt_support}
    The Schmidt rank $r_a$ of the nonzero loop bond vector is positive.
\end{lemma}

\begin{proof}\leanok
    If $r_a=0$, then the intermediate index set in the factorization
    $\Phi_a=X_aY_a$ is empty. Every entry of $X_aY_a$ is therefore zero,
    contrary to $\varphi_a\ne0$.
\end{proof}

\begin{theorem}[Injectivity of the minimal Beigi loop tensor]
    \label{thm:beigi_minimal_loop_tensor_injective}
    \lean{MPSTensor.BeigiSectorGraphData.minimalLoopCoordinateTensor_isInjective}
    \lean{MPSTensor.BeigiSectorGraphData.minimalLoopTensor_isInjective}
    \lean{MPSTensor.BeigiSectorGraphData.minimalLoopTensor_isNormal}
    \leanok
    \uses{def:beigi_loop_schmidt_support, def:injective, def:normal}
    The tensor $\widetilde C_a$ on the Schmidt support is injective, as is its
    image under the one-site unitary. In particular, the physical tensor is
    normal.
\end{theorem}

\begin{proof}\leanok
    In a rank factorization $\Phi_a=X_aY_a$ through
    $r_a=\rank\Phi_a$, the columns of $Y_a$ and the rows of $X_a$ each span
    $\C^{r_a}$. Their pairwise outer products therefore span
    $M_{r_a}(\C)$, and these outer products are precisely the loop-sector
    letters of $\widetilde C_a$. A unitary change of physical coordinates
    preserves this span. One-site injectivity implies normality.
\end{proof}

\begin{theorem}[Transfer fixed point of a normalized Beigi loop tensor]
    \label{thm:beigi_normalized_minimal_loop_tensor_rfp}
    \lean{MPSTensor.BeigiSectorGraphData.loopBondNorm_pos}
    \lean{MPSTensor.BeigiSectorGraphData.loopBondNorm_ne_zero}
    \lean{MPSTensor.BeigiSectorGraphData.loopBondNorm_eq_frobeniusNorm}
    \lean{MPSTensor.BeigiSectorGraphData.transferMap_minimalLoopCoordinateTensor}
    \lean{MPSTensor.BeigiSectorGraphData.transferMap_minimalLoopCoordinateTensor_comp_self}
    \lean{MPSTensor.BeigiSectorGraphData.transferMap_normalizedMinimalLoopCoordinateTensor}
    \lean{MPSTensor.BeigiSectorGraphData.normalizedMinimalLoopCoordinateTensor_isTransferIdempotent}
    \lean{MPSTensor.BeigiSectorGraphData.normalizedMinimalLoopTensor_isTransferIdempotent}
    \lean{MPSTensor.BeigiSectorGraphData.normalizedMinimalLoopTensor_isInjective}
    \lean{MPSTensor.BeigiSectorGraphData.normalizedMinimalLoopTensor_isNormal}
    \leanok
    \uses{def:beigi_normalized_minimal_loop_tensor,
        thm:beigi_minimal_loop_tensor_injective,
        def:mps_transfer_idempotence}
    The transfer map of the unnormalized coordinate tensor is
    \begin{align}
        \mathcal E_a(Z)
         & =
        \tr(X_a^*X_aZ)Y_aY_a^*, \notag \\
        \mathcal E_a^2
         & =
        \nu_a^2\mathcal E_a. \notag
    \end{align}
    Hence the transfer map of $\widehat C_a$ is idempotent. This remains
    true after the physical unitary rotation. The normalized physical tensor
    is injective and therefore normal.
\end{theorem}

\begin{proof}\leanok
    Summing the Kraus terms over the loop-sector coordinates gives the stated
    rank-one formula. Cyclicity of the trace and $\Phi_a=X_aY_a$ give
    $\tr(X_a^*X_aY_aY_a^*)=\tr(\Phi_a^*\Phi_a)=\nu_a^2$. Scaling every
    letter by $\nu_a^{-1}$ scales the transfer map by $\nu_a^{-2}$, which
    makes it idempotent. A unitary change of physical coordinates leaves the
    transfer map invariant, and the nonzero scalar preserves injectivity.
\end{proof}

\begin{theorem}[Periodic vector of a positive Beigi loop]
    \label{thm:beigi_loop_tensor_periodic_vector}
    \lean{MPSTensor.BeigiSectorGraphData.mpv_loopCoordinateTensor}
    \lean{MPSTensor.BeigiSectorGraphData.mpv_loopTensor}
    \lean{MPSTensor.BeigiSectorGraphData.mpv_minimalLoopCoordinateTensor}
    \lean{MPSTensor.BeigiSectorGraphData.mpv_minimalLoopTensor}
    \lean{MPSTensor.BeigiSectorGraphData.mpv_normalizedMinimalLoopTensor}
    \leanok
    \uses{def:beigi_loop_product_state, def:beigi_loop_schmidt_support,
        def:beigi_normalized_minimal_loop_tensor, lem:mpv_rotate_physical}
    For every $N>0$, the periodic matrix product vector of $C_a$ has
    coefficient $\prod_{n=0}^{N-1}\varphi_a(r_n,l_{n+1})$ on a cyclic
    configuration in sector~$a$, and coefficient zero on every other sector
    configuration. Consequently, the periodic vector of the physical loop
    tensor is exactly the sitewise-unitary image of this cyclic
    product-of-pairs state. The same identities hold for the tensor on the
    minimal Schmidt support. For the normalized minimal tensor it is
    $\nu_a^{-N}\Psi_a^{(N)}$.
\end{theorem}

\begin{proof}\leanok
    Write $D_a^{\mathrm L}$ for the left-factor dimension of sector~$a$.
    Expanding the closed matrix product on a sector configuration
    $((r_n,l_n))_{n=0}^{N-1}$ gives
    \begin{align}
         & \sum_{g\in\{0,\ldots,D_a^{\mathrm L}-1\}^N}
        \prod_{n=0}^{N-1}
        \delta_{g_n,l_n}
        \varphi_a\!\left(r_n,g_{n+1}\right) =
        \prod_{n=0}^{N-1}\varphi_a(r_n,l_{n+1}), \notag
    \end{align}
    where all site labels are read cyclically. Indeed, the Kronecker deltas
    leave only the assignment $g_n=l_n$.

    For the minimal tensor, let $S^{(r,l)}$ be the rectangular matrix whose
    only nonzero row is row~$l$, equal there to row~$r$ of $X_a$. Then
    $C_a^{(r,l)}=S^{(r,l)}Y_a$ and
    $\widetilde C_a^{(r,l)}=Y_aS^{(r,l)}$. For every nonempty word there is
    therefore a rectangular matrix $R$ such that its two products are $RY_a$
    and $Y_aR$, respectively. Cyclicity of the trace gives
    $\tr(RY_a)=\tr(Y_aR)$, proving equality of all positive-length periodic
    vectors.

    If $A_a$ denotes the physical loop tensor and $U$ the one-site unitary,
    then for a physical configuration $s=(s_0,\ldots,s_{N-1})$ one has
    \begin{align}
        \tr(A_a^{s_0}\cdots A_a^{s_{N-1}})
         & =
        \sum_{t\in\{0,\ldots,d-1\}^N}
        \left(\prod_{n=0}^{N-1}U_{s_n,t_n}\right)
        \tr(C_a^{t_0}\cdots C_a^{t_{N-1}}). \notag
    \end{align}
    Substitution of the preceding cyclic coefficient formula gives the stated
    sitewise-unitary image.
\end{proof}

\begin{theorem}[Spatial conjugation of complementary interaction products]
    \label{thm:beigi_complement_product_conjugation}
    \lean{MPSTensor.BeigiSectorGraphData.singleKrausMap_groundBondProduct}
    \leanok
    \uses{def:beigi_sector_graph}
    Suppose that $A$ is equipped with Beigi sector data, with one-site spatial
    unitary $U$. For $N\geq2$, let $Q_N$ be the product of the translated
    complementary parent interactions and let $\widetilde Q_N$ be the
    corresponding product in sector coordinates. Then
    $(U^*)^{\otimes N}Q_NU^{\otimes N} = \widetilde Q_N$.
\end{theorem}

\begin{proof}\leanok
    The two-site coordinate identity gives
    $(U^*\otimes U^*)(\Id-h)(U\otimes U)=\widetilde Q$ for one bond.
    Conjugation by $(U^*)^{\otimes N}$ carries every translated copy of this
    bond identity to the corresponding translated sector interaction. Since
    unitary conjugation distributes over an ordered product,
    \begin{align}
        (U^*)^{\otimes N}
        \left(\prod_{n=0}^{N-1}Q_n\right)
        U^{\otimes N}
         & =
        \prod_{n=0}^{N-1}
        \left((U^*)^{\otimes N}Q_nU^{\otimes N}\right) =
        \prod_{n=0}^{N-1}\widetilde Q_n =
        \widetilde Q_N, \notag
    \end{align}
    where $Q_n$ and $\widetilde Q_n$ are the translated physical and sector
    complementary interactions, respectively.
\end{proof}

\begin{theorem}[Fixed complementary-product vectors have zero parent energy]
    \label{thm:beigi_complement_product_fixed_vector}
    \lean{MPSTensor.BeigiSectorGraphData.mem_ker_parentHamiltonian_of_groundBondProduct_mulVec_eq_self}
    \leanok
    \uses{def:beigi_sector_graph}
    Suppose that $A$ is equipped with Beigi sector data. For $N\geq2$, let
    $h_0,\ldots,h_{N-1}$ be its translated length-two parent interactions and
    set $Q_N=\prod_{n=0}^{N-1}(\Id-h_n)$. Every vector fixed by $Q_N$
    belongs to the kernel of the periodic parent Hamiltonian
    $H_N=\sum_n h_n$.
\end{theorem}

\begin{proof}\leanok
    Under the Euclidean identification $E$, the matrix product of the
    complementary interactions is carried to
    \begin{align}
        Q_N^{\mathrm{E}}
         & =
        E^{-1}Q_NE =
        \prod_{n=0}^{N-1}(\Id-h_n^{\mathrm{E}}). \notag
    \end{align}
    The operators $h_n^{\mathrm{E}}$ are orthogonal projections. Beigi sector
    data supply pairwise commutativity, which Euclidean transport preserves.
    For all $0\leq m,n<N$, one has
    $h_m^{\mathrm{E}}h_n^{\mathrm{E}}
        =h_n^{\mathrm{E}}h_m^{\mathrm{E}}$. For any commuting orthogonal
    projections $P_0,\ldots,P_{N-1}$,
    \begin{align}
        \prod_{n=0}^{N-1}(\Id-P_n)
         & =
        \operatorname{proj}_{\bigcap_{n=0}^{N-1}\ker P_n}. \notag
    \end{align}
    Applying this identity to the $h_n^{\mathrm{E}}$ gives
    \begin{align}
        Q_N^{\mathrm{E}}
         & =
        \operatorname{proj}_{\bigcap_{n=0}^{N-1}\ker h_n^{\mathrm{E}}},
        \notag \\
        \bigcap_{n=0}^{N-1}\ker h_n^{\mathrm{E}}
         & =
        \ker\!\left(\sum_{n=0}^{N-1}h_n^{\mathrm{E}}\right) =
        \ker H_N^{\mathrm{E}}. \notag
    \end{align}
    Thus $Q_Nv=v$ implies $E^{-1}v\in\ker H_N^{\mathrm{E}}$, and conjugating
    back by $E$ gives $v\in\ker H_N$.
\end{proof}

\begin{theorem}[A positive Beigi loop gives a parent ground state]
    \label{thm:beigi_loop_product_state_ground_state}
    \lean{MPSTensor.BeigiSectorGraphData.loopProductState_mem_ker_parentHamiltonian}
    \lean{MPSTensor.BeigiSectorGraphData.loopProductState_mem_parentHamiltonianGroundSpaceES}
    \leanok
    \uses{def:beigi_sector_graph, def:beigi_loop_product_state,
        thm:beigi_complement_product_conjugation,
        thm:beigi_complement_product_fixed_vector}
    Let $a\to a$ be a positive loop. For every $N\geq2$, its physical
    cyclic product-of-pairs state belongs to the kernel of the length-two
    periodic parent Hamiltonian. Equivalently, its Euclidean representative
    belongs to the parent ground space. This is only a membership assertion;
    it does not assert that the loop states span the ground space.
\end{theorem}

\begin{proof}\leanok
    In sector coordinates, the complete product $\widetilde Q_N$ of the
    complementary two-site interactions is block diagonal over cyclic sector
    sequences. Write $\Phi_a(k;x)$ for the coefficient of the loop cyclic
    product in the block indexed by $k$. Its support lies in sector~$a$, so
    for every $k\ne(a,\ldots,a)$ and every edge-index tuple $x$, one has
    $\Phi_a(k;x)=0$. Thus only the constant cyclic sector sequence
    contributes. On this block, $\widetilde Q_N$ is the tensor product of the
    edge projectors. Since $\varphi_a\in\ran P_{a,a}$,
    \begin{align}
        \widetilde Q_N\Phi_a
         & =
        P_{a,a}^{\otimes N}\Phi_a =
        \Phi_a, \notag \\
        \Phi_a(t_0,\ldots,t_{N-1})
         & =
        \prod_{n=0}^{N-1}\varphi_a(r_n,l_{n+1}). \notag
    \end{align}
    If $U$ is the one-site spatial unitary and
    $\Psi_a=U^{\otimes N}\Phi_a$ is the physical loop state, unitary
    conjugation gives
    \begin{align}
        Q_N\Psi_a
         & =
        U^{\otimes N}\widetilde Q_N\Phi_a =
        U^{\otimes N}\Phi_a =
        \Psi_a. \notag
    \end{align}
    Theorem~\ref{thm:beigi_complement_product_fixed_vector} now yields
    $\Psi_a\in\ker H_N$.
\end{proof}

\begin{theorem}[Positive Beigi loops span the parent ground space]
    \label{thm:beigi_loop_product_states_span_ground_space}
    \lean{MPSTensor.BeigiSectorGraphData.loopCyclicProduct_ne_zero}
    \lean{MPSTensor.BeigiSectorGraphData.loopCyclicProduct_inner_eq_zero_of_ne}
    \lean{MPSTensor.BeigiSectorGraphData.loopProductState_ne_zero}
    \lean{MPSTensor.BeigiSectorGraphData.loopProductStateES_ne_zero}
    \lean{MPSTensor.BeigiSectorGraphData.loopProductStateES_inner_eq_zero_of_ne}
    \lean{MPSTensor.BeigiSectorGraphData.loopProductStateES_linearIndependent}
    \lean{MPSTensor.BeigiSectorGraphData.span_loopProductStateES_eq_parentHamiltonianGroundSpaceES}
    \lean{MPSTensor.BeigiSectorGraphData.span_loopProductStateES_eq_parentHamiltonianGroundSpaceES_of_hasBNTSectorData}
    \leanok
    \uses{thm:beigi_loop_product_state_ground_state,
        cor:beigi_parent_ground_space_loop_number}
    Suppose that the parent-ground-space dimension is eventually constant.
    For every $N>2$, the positive-loop product states form a nonzero pairwise
    orthogonal family and
    \begin{align}
        \spn\{\Psi_a^{(N)}:a\to a\text{ is a positive loop}\}
         & =
        \ker H_N. \notag
    \end{align}
    In particular, the same conclusion holds under the BNT and all-chain
    nearest-neighbor parent-ground-space hypotheses.
\end{theorem}

\begin{proof}\leanok
    Write $\Phi_a^{(N)}$ for the cyclic product in sector coordinates and
    $\Psi_a^{(N)}=U^{\otimes N}\Phi_a^{(N)}$ for the physical state, where
    $U$ is the one-site spatial unitary. Choose $q$ with
    $\varphi_a(q)\ne0$. On the constant edge-index configuration,
    \begin{align}
        \Phi_a^{(N)}(q,\ldots,q)
         & =
        \prod_{n=0}^{N-1}\varphi_a(q) \notag \\
         & \ne
        0. \notag
    \end{align}
    Thus $\Phi_a^{(N)}\ne0$, and unitarity gives
    $\Psi_a^{(N)}=U^{\otimes N}\Phi_a^{(N)}\ne0$. If $a\ne b$, then every
    sector configuration $s$ satisfies $\Phi_a^{(N)}(s)=0$ or
    $\Phi_b^{(N)}(s)=0$, and hence
    \begin{align}
        \langle\Phi_a^{(N)},\Phi_b^{(N)}\rangle
         & =
        \sum_s\overline{\Phi_a^{(N)}(s)} \Phi_b^{(N)}(s) =
        0. \notag
    \end{align}
    Moreover,
    \begin{align}
        \langle\Psi_a^{(N)},\Psi_b^{(N)}\rangle
         & =
        \langle U^{\otimes N}\Phi_a^{(N)},
        U^{\otimes N}\Phi_b^{(N)}\rangle =
        \langle\Phi_a^{(N)},\Phi_b^{(N)}\rangle =
        0. \notag
    \end{align}
    Thus the loop states are linearly independent. Their span is contained in
    $\ker H_N$ by
    Theorem~\ref{thm:beigi_loop_product_state_ground_state}. Its dimension is
    the number of positive loops, which equals $\dim\ker H_N$ by
    Corollary~\ref{cor:beigi_parent_ground_space_loop_number}. The inclusion
    is consequently an equality.
\end{proof}

\begin{theorem}[All-chain nearest-neighbor commuting parent-Hamiltonian ground spaces imply a renormalization fixed point]
    \label{thm:nncph_implies_rfp}
    \lean{MPSTensor.nncph_implies_rfp}
    \leanok
    \uses{def:mps_transfer_idempotence, def:bnt, def:has_nncph_ground_spaces,
        def:normal, thm:beigi_loop_product_states_span_ground_space,
        thm:beigi_normalized_minimal_loop_tensor_rfp}
    Let $B$ be normal and in left-canonical form. Suppose that the family
    $A_j$ is a BNT for $B$, and that the nearest-neighbor parent Hamiltonian
    of $B$ has the all-chain ground-space property of
    Definition~\ref{def:has_nncph_ground_spaces}: for every $N>2$, the
    length-two translated parent interactions commute,
    $h_i(B,2)h_j(B,2)=h_j(B,2)h_i(B,2)$, the vector $V^{(N)}(B)$ has zero
    energy, and
    $\ker H_2^{(N)}(B) = \spn\{V^{(N)}(A_j):j=1,\ldots,g\}$.
    Then $B$ is a renormalization fixed point.

    This is the normal left-canonical representative case of the reverse
    implication in Theorem~3.10 of \cite{Cirac2016MPDO_arXiv}. The source
    theorem is stated for a tensor in canonical form; the unrestricted
    canonical-form statement is not asserted here.
\end{theorem}

\begin{proof}
    \leanok
    The commuting interaction determines Beigi's finite sector graph. Its
    positive loops span the parent ground space, and the normalized minimal
    loop tensors $L_j$ are normal renormalization fixed points. For the normal
    representative $B$,
    \begin{align}
        \forall^\infty N,\quad
        V^{(N)}(B)
         & \in
        \spn\{V^{(N)}(L_j):j=1,\ldots,g\}. \notag
    \end{align}
    Normal-tensor orthogonality applied to this eventual-span relation first
    identifies the bond dimensions and then gives an index $j$, an invertible
    matrix $X$, and a scalar $\zeta$ with $|\zeta|=1$ such that
    $L_j^i = \zeta X B^i X^{-1}$.
    Transfer idempotence is preserved by these transformations.
\end{proof}

\begin{lemma}[The unit-weight basis direct sum is a BNT]
    \label{lem:bnt_basis_direct_sum}
    \lean{MPSTensor.IsBNTCanonicalForm.isBNT_basisDirectSum}
    \leanok
    Let $A_1,\ldots,A_g$ be the distinct basis tensors of a BNT canonical
    form. Then their multiplicity-one, unit-weight direct sum
    \[
        B=\bigoplus_{j=1}^g A_j
    \]
    has the same family $A_1,\ldots,A_g$ as an algebraic BNT.
\end{lemma}

\begin{proof}
    \leanok
    Each $A_j$ is normal and left-canonical. The direct-sum state is the sum
    of the states of the $A_j$ with coefficient one, and the eventual linear
    independence is part of the given BNT canonical form.
\end{proof}

\begin{lemma}[Common physical unitaries preserve mixed transfer maps]
    \label{lem:mixed_transfer_common_physical_unitary}
    \lean{MPSTensor.mixedTransferMap₂_rotatePhysical}
    \leanok
    Let $A=(A_i)_i$ and $B=(B_i)_i$ be two tensors, possibly with
    different bond dimensions, and let $u$ be unitary on the physical index.
    If
    \[
        \widetilde A_i=\sum_j u_{ij}A_j,
        \qquad
        \widetilde B_i=\sum_j u_{ij}B_j,
    \]
    then
    \[
        \mathcal E_{\widetilde A,\widetilde B}
        =
        \mathcal E_{A,B}.
    \]
\end{lemma}

\begin{proof}
    \leanok
    Expand both physical sums. The coefficient of
    $A_jX B_k^\dagger$ is
    $\sum_i u_{ij}\overline{u_{ik}}=\delta_{jk}$, so only the diagonal
    terms remain.
\end{proof}

\begin{lemma}[Gauge-phase transport of mixed-transfer vanishing]
    \label{lem:mixed_transfer_gauge_phase_zero}
    \lean{MPSTensor.mixedTransferMap₂_eq_zero_of_gaugePhaseEquiv}
    \lean{MPSTensor.mixedTransferMap₂_cast_eq_zero_iff}
    \leanok
    \uses{def:gauge_phase_equiv}
    Let $U,A$ and $V,B$ be pairs of tensors with respective common bond
    dimensions. Suppose that $\zeta_A,\zeta_B\ne 0$ and that $X_A,X_B$ are
    invertible matrices such that, for every physical index $i$,
    \begin{align}
        A^i &= \zeta_A X_A U^i X_A^{-1}, \\
        B^i &= \zeta_B X_B V^i X_B^{-1}.
    \end{align}
    Then
    \begin{align}
        \mathcal E_{A,B}=0
        \quad\Longrightarrow\quad
        \mathcal E_{U,V}=0.
    \end{align}
    Identifying equal bond dimensions on either tensor leg preserves this
    vanishing condition in both directions.
\end{lemma}

\begin{proof}
    \leanok
    Set $D_A=\zeta_A I$ and $D_B=\zeta_B I$. For every rectangular matrix
    $Z$, covariance of the mixed transfer map gives
    \begin{align}
        \mathcal E_{A,B}(X_A Z X_B^\dagger)
        &=
        X_A D_A\,\mathcal E_{U,V}(Z)\,(X_B D_B)^\dagger .
    \end{align}
    If $\mathcal E_{A,B}=0$, the left-hand side vanishes. The two outer
    factors are invertible, so cancellation gives
    $\mathcal E_{U,V}(Z)=0$ for every $Z$.
\end{proof}

\begin{lemma}[Equal bond-dimension identifications preserve mixed-transfer vanishing]
    \label{lem:mixed_transfer_dimension_identification_zero}
    \lean{MPSTensor.mixedTransferMap₂_cast_eq_zero_iff}
    \leanok
    Identifying either bond dimension with an equal natural number does not
    change whether the rectangular mixed transfer map vanishes.
\end{lemma}

\begin{proof}
    \leanok
    After the two dimension identifications, both tensors and their mixed
    transfer map are unchanged.
\end{proof}

\begin{lemma}[Distinct normalized Beigi loops are locally orthogonal]
    \label{lem:beigi_normalized_minimal_loop_mixed_transfer_zero}
    \lean{MPSTensor.BeigiSectorGraphData.normalizedMinimalLoopTensor_mixedTransferMap₂_eq_zero_of_ne}
    \leanok
    \uses{lem:mixed_transfer_common_physical_unitary}
    If $a$ and $b$ are distinct positive loops in the sector graph, then the
    mixed transfer map of their normalized minimal loop tensors vanishes:
    \[
        \mathcal E_{L_a,L_b}=0.
    \]
\end{lemma}

\begin{proof}
    \leanok
    In sector coordinates, every physical letter of $L_a$ is supported in
    sector $a$, while every physical letter of $L_b$ is supported in sector
    $b$. Hence, for every rectangular matrix $Z$ and every physical index $i$,
    \begin{align}
        L_a^i Z (L_b^i)^\dagger &= 0, \\
        \mathcal E_{L_a,L_b}(Z)
        &= \sum_i L_a^i Z (L_b^i)^\dagger = 0.
    \end{align}
    The common one-site physical unitary preserves the mixed transfer map by
    Lemma~\ref{lem:mixed_transfer_common_physical_unitary}.
\end{proof}

\begin{theorem}[Multi-sector NNCPH ground spaces imply a fixed point for the basis direct sum]
    \label{thm:nncph_implies_rfp_basis_direct_sum}
    \lean{MPSTensor.IsBNTCanonicalForm.isTransferIdempotent_basisDirectSum_of_hasNNCPHGroundSpaces}
    \leanok
    \uses{def:mps_transfer_idempotence, def:has_nncph_ground_spaces,
        lem:bnt_basis_direct_sum,
        thm:beigi_loop_product_states_span_ground_space,
        thm:beigi_normalized_minimal_loop_tensor_rfp,
        lem:beigi_normalized_minimal_loop_mixed_transfer_zero,
        lem:mixed_transfer_gauge_phase_zero,
        lem:mixed_transfer_dimension_identification_zero,
        thm:cpsv_normal_mpv_phase_gauge,
        thm:direct_sum_transfer_idempotent_iff_pairwise_mixed}
    Let $A_1,\ldots,A_g$ be the distinct basis tensors of a BNT canonical
    form, and set
    \[
        B=\bigoplus_{j=1}^g A_j .
    \]
    Suppose that, for every $N>2$, the translated length-two parent
    interactions of $B$ commute, every $V^{(N)}(A_j)$ has zero energy, and
    \[
        \ker H_2^{(N)}(B)
        =
        \spn\{V^{(N)}(A_j):j=1,\ldots,g\}.
    \]
    Then $B$ is a renormalization fixed point.

    This is the multiplicity-one, unit-weight representative case of the
    reverse implication in Theorem~3.10 of
    \cite{Cirac2016MPDO_arXiv}. It does not assert the printed theorem for
    repeated copies or arbitrary raw weights.
\end{theorem}

\begin{proof}
    \leanok
    Beigi's positive-loop product states span each parent ground space.
    Consequently every normal basis tensor $A_j$ agrees, up to a unit-modulus
    scalar and an invertible virtual change of basis, with a normalized
    minimal loop tensor $L_{\ell(j)}$. Distinctness of the BNT basis makes
    $j\mapsto\ell(j)$ injective. Thus there are $\zeta_j\ne 0$ and invertible
    $X_j$ such that
    \begin{align}
        L_{\ell(j)}^i &= \zeta_j X_j A_j^i X_j^{-1}.
    \end{align}

    Each $L_{\ell(j)}$ has an idempotent transfer map. If $j\ne k$, then
    $L_{\ell(j)}$ and $L_{\ell(k)}$ have disjoint sector support, so their
    mixed transfer map vanishes. Lemma~\ref{lem:mixed_transfer_gauge_phase_zero}
    transports this vanishing through the two virtual changes of basis:
    \begin{align}
        \mathcal E_{L_{\ell(j)},L_{\ell(k)}}=0
        \quad\Longrightarrow\quad
        \mathcal E_{A_j,A_k}=0.
    \end{align}
    Thus the diagonal mixed transfer maps of the family $A_j$ are idempotent
    and the off-diagonal ones vanish. The blockwise criterion is
    \begin{align}
        \mathcal E_B\circ\mathcal E_B=\mathcal E_B
        \quad\Longleftrightarrow\quad
        \mathcal E_{A_j,A_k}\circ\mathcal E_{A_j,A_k}
        =\mathcal E_{A_j,A_k}
    \end{align}
    for every pair $j,k$, and proves the claim.
\end{proof}
